% =====================================================================
%  DRAFT expansion of Section 2, written after reading all five papers.
%
%  VERIFY EVERY CHARACTERISATION AGAINST THE PAPER BEFORE SUBMITTING.
%  I have read them and the claims below are sourced, but you are the
%  author and you should be able to defend each sentence.
%
%  Replaces 2.2, extends 2.3, and adds 2.5 and 2.6.
%  The fifth column of Table 2.1 is YOUR argument - none of these papers
%  asks what a quality filter removes. That is what makes the table worth
%  printing rather than listing.
% =====================================================================

\subsection{Taxonomies of disagreement}

Asking what disagreement \emph{is} produces more than one answer, and the answers do not reduce
to a single decomposition. Four accounts are relevant here, and they differ in where they locate
the disagreement rather than merely in vocabulary.

Lindahl (2024) distinguishes boundary, label, and existence disagreement: annotators may differ
about where the task's boundary falls, about which label applies within it, or about whether a
valid answer exists at all. Jiang and de Marneffe (2022) work bottom-up from data rather than
top-down from theory, developing a taxonomy of ten disagreement sources over MNLI and ChaosNLI
items, grouped into three high-level classes. \emph{Uncertainty in sentence meaning} covers
lexical ambiguity, implicature, presupposition, probabilistic enrichment, and imperfection;
\emph{underspecification in guidelines} covers coreference, temporal reference, and interrogative
hypotheses; \emph{annotator behaviour} covers accommodating minimally added content and high
lexical overlap heuristics. The three classes locate disagreement in the item, in the task
definition, and in the annotator respectively.

Weber-Genzel et al.\ (2024) draw a different line, between human label variation --- different
labels assigned for valid reasons --- and annotation error, where a label is assigned for an
invalid reason. Their contribution is methodological as much as conceptual: in VariErr NLI they
operationalise the distinction with a two-round procedure in which annotators first explain each
label and then judge the validity of others' label--explanation pairs, producing 7,732 validity
judgments over 1,933 explanations for 500 re-annotated MNLI items. The distinction is therefore
carried by the \emph{reason}, not by the label. Sang and Stanton (2022) locate disagreement in
the annotator instead, developing a multidimensional scale from interviews and concept mapping
and testing it with 170 hate-speech annotators, and reporting that facets of individual
difference such as age and personality relate to labelling decisions.

\begin{table}[htbp]
\centering
\caption{Four accounts of disagreement, and what each implies about rater filtering.}
\label{tab:taxonomies}
\small
\begin{tabularx}{\textwidth}{p{2.5cm}p{2.2cm}Xp{2.6cm}}
\toprule
\textbf{Source} & \textbf{Domain} & \textbf{Categories} & \textbf{Located in} \\
\midrule
Lindahl (2024) & argumentation & boundary / label / existence & task, answer, or whether an
answer exists \\
Jiang \& de Marneffe (2022) & NLI & 10 sources in 3 classes: uncertainty in sentence meaning;
underspecification in guidelines; annotator behaviour & the item, the guidelines, or the
annotator \\
Weber-Genzel et al.\ (2024) & NLI & human label variation vs annotation error, separated by
validity of the stated reason & the reason behind the label \\
Sang \& Stanton (2022) & hate speech & multidimensional scale of annotator individual
differences & stable annotator characteristics \\
\bottomrule
\end{tabularx}
\end{table}

These accounts matter for the present study in a specific way, and the point is not one any of
them makes. A rater quality filter is an operation that removes annotators. Read against
Jiang and de Marneffe's taxonomy, such a filter is aimed at their third class --- annotator
behaviour --- while their first two classes are precisely the signal a disagreement-aware
analysis wants to keep. A filter that cannot distinguish the classes removes both. Read against
Weber-Genzel et al., the difficulty sharpens: separating error from variation required them to
collect explanations and validity judgments, because the distinction lives in the reason and not
in the label pattern. The DICES filter, like most published quality filters, records no reasons
at all. The error/variation distinction is therefore not merely unmeasured in this dataset but
unavailable in principle, which is why Section~\ref{sec:published-filter} treats the published
exclusion list as a single fixed operation rather than a criterion that could be re-derived.

Sang and Stanton's finding adds a further caution. If disagreement is partly a function of stable
annotator characteristics, then filtering on behavioural signals risks removing \emph{people}
rather than \emph{errors} --- which is testable, and is what RQ1 tests.

This also fixes the scope of what this paper measures. The analysis operates on categorical
item-level labels: soft-label vectors over No, Yes, and Unsure, compared with TVD and JSD.
Under Lindahl's decomposition this captures label disagreement directly but cannot separate
boundary from existence disagreement, and under Weber-Genzel et al.'s it cannot separate
variation from error at all. Two raters who disagree because they read the task differently are
indistinguishable here from two who read it the same way and judged differently. The measurements
therefore quantify one observable form of disagreement, and the negative persona result in
particular is scoped to that form.

\subsection{What to do with disagreement once it has been measured}
\label{sec:whattodo}

Measuring disagreement and deciding what to do with it are separable problems, and the present
study addresses only the first. Two lines of work address the second and are worth stating
because they define what a measurement like this one is \emph{for}.

Gordon et al.\ (2022) pose the question directly: whose labels should a model learn to emulate?
Their answer, jury learning, replaces implicit majority vote with an explicit choice of which
people or groups determine a prediction, and in what proportion. Architecturally they model every
annotator in a dataset, sample from those models to populate a jury, then run inference. The
framing is directly relevant here: a between-group spread measurement is a precondition for jury
learning rather than a competitor to it. If a quality filter has silently removed one demographic
axis from a corpus, a jury drawn from that corpus inherits the removal, and no amount of care in
composing the jury will recover what is no longer there.

Romberg (2022) takes a different route to the same concern. Rather than changing whose labels
count, PerspectifyMe enriches a prediction with a subjectivity score derived from the annotation
process, so that a model reports both an aggregated hard label and an indication of which items
elicited genuinely different views. Applied to argument concreteness, this treats contested items
as something to flag rather than to resolve --- an alternative to removing raters that leaves the
disagreement visible downstream.

\subsection{Summary of the gap}

Filtering and persona-prompting have each been examined, and disagreement has been taxonomised
(Lindahl, 2024; Jiang \& de Marneffe, 2022; Weber-Genzel et al., 2024), traced to annotator
characteristics (Sang \& Stanton, 2022), and modelled downstream (Gordon et al., 2022;
Romberg, 2022). What has not been done is to measure both operations on the same items, with the
same metric, against the same null, and to ask whether they converge on the same reduced view.
That is the question this paper addresses.
